\chapter*{Workshop 7: Ship a reproducible example}\label{workshop-07}
\addcontentsline{toc}{chapter}{W7: Ship a reproducible example}

\section{Pin the thing a reader can actually run}

The companion lives in its own public repository: \url{https://github.com/flyingrobots/bijou-job-monitor}. This edition uses tag \texttt{book-v1} and the exact commit recorded on the version page. The package lock pins published Bijou 7.2.0 and the supporting toolchain. The archived records elsewhere in the book discuss several versions and aspirations; they are not the dependency manifest for this application.

\subsection{Reproduce a checkpoint}

Clone the repository, check out the recorded commit or edition tag, and run \texttt{npm ci}. Run the checkpoint command for the workshop being read, then its checks. The first two checkpoints display a model and its modes. The third opens the complete interactive application. The fourth reproduces stale-result rejection, and the fifth exercises narrow layout and expanded text. The sixth workshop runs the test evidence; this workshop inspects packaging. The checkpoints share the finished modules instead of maintaining five independent copies that can drift.

This is a reproducible route through the construction, not a claim that every checkpoint is a separate historical release. When making an incompatible change to an example, create another edition tag and update the book's recorded reference. A reader should be able to return to the code that produced the explanation.

\subsection{Verify the distributed source}

Run \texttt{npm pack --dry-run} and inspect the listed files. The code package includes the full Apache 2.0 license, its NOTICE, source-file copyright statements, checkpoints and exercises. Git checkout also retains the package lock. Dependencies retain their own applicable notices if they are redistributed with a package. The copyright page in the book does not replace these files.

\section{Keep claims attached to their environment}

\subsection{An environment record has limits}

The version page identifies Node, npm, the companion commit, the published Bijou release and the locally tested host. The companion's \path{VERIFICATION.json} records fifteen passing tests and the pseudo-terminal observations. GitHub Actions runs the pinned checks on Linux; its current result belongs to that run and should be read from the repository. No human screen-reader session or broad terminal-emulator matrix was performed for this edition. Those are open acceptance questions rather than implied guarantees.

\subsection{Make corrections findable}

Use the companion's issue tracker for book-example errata. Include the workshop, edition commit, command, environment, expected result and observed result. A correction becomes more useful when another reader can reproduce it. Keep the original edition available while a later edition explains the repair and its compatibility effect.

\section{Let the book be a working specimen}

The monitor already lets readers start jobs, switch to logs, change selection, filter failures and resize the view. The explicit CLI modes let them inspect the same fixture as static text, TSV or accessible prose. Its tests expose state and cells without requiring an image. The printed explanations remain complete; the running application supplies another way to examine their claims.

A future inspector could make model snapshots and layout rectangles visible inside the application. That capability is not implemented in this checkpoint. It should be added when it answers a concrete learning question and arrives with its own evidence. The useful flourish is the connection between the explanation and a working surface, not an unverified menu promising every debugging feature.
